\section{11 February 2022: The Decree, Read Word by Word}

\subsection{The audience of 4 February}

The Fraternity's account of what produced the decree is given in its
communiqué of 21 February 2022. On Friday 4 February 2022 two of its
members---Father Benoît Paul-Joseph, superior of the district of France,
and Father Vincent Ribeton, rector of the seminary at Wigratzbad---were
received in private audience by Francis for nearly an hour. In the
Fraternity's telling the pope, hearing of the origins of 1988, said that
the founders' gesture ``should be `preserved, protected and
encouraged'\,''; and ``the Pope made it clear that institutes such as the
Fraternity of St.\ Peter are not affected by the general provisions of the
Motu Proprio \work{Traditionis Custodes}, since the use of the ancient
liturgical books was at the origin of their existence and is provided for
in their constitutions.'' The Holy Father ``subsequently sent a decree
signed by him and dated February 11.''\endnote{``Official communiqué of
the Priestly Fraternity of St.\ Peter'' (fssp.org), Fribourg, 21 February
2022, read 2026-07-25. A party account of a private audience; the words
attributed to the pope are the Fraternity's report of what was said, and
no Holy See record of the audience was located. This account reports the
report and asserts nothing about what was said. The date 11 February is
independently established by the decree itself.}

That report should be read with care, because the decree does not say
what the report says the pope said. The report states a
principle---institutes of this kind are not affected by the general
provisions of \work{Traditionis custodes}---which would resolve the whole
question of art.~8. The decree states no such principle. It grants a
faculty.

\subsection{The Latin}

The decree is short enough to set out in full, and the Latin matters
because the differences between it and the English rendering the
Fraternity publishes are not trivial.

\begin{quote}\small
\latin{Sanctus Pater Franciscus, omnibus et singulis sodalibus Instituti
vitae consecratae ``Fraternitas Sancti Petri'' nuncupati, die 18 iulii
1988 erecti et a Sancta Sede pontificii iuris declarati, facultatem
concedit celebrandi sacrificium Missae, sacramentorum necnon alios sacros
ritus, sicut et persolvendi Officium divinum, iuxta editiones typicas
librorum liturgicorum, scilicet Missalis, Ritualis, Pontificalis et
Breviarii, anno 1962 vigentium.}

\latin{Qua facultate uti poterunt in ecclesiis vel oratoriis propriis,
alibi vero nonnisi de consensu Ordinarii loci, excepta Missae privatae
celebratione.}

\latin{Quibus rite servatis, Sanctus Pater etiam suadet ut sedulo
cogitetur, quantum fieri potest, de statutis in litteris apostolicis motu
proprio datis Traditionis Custodes.}

\latin{Datum Romae, Sancti Petri, die XI mensis Februarii, in memoria
Beatae Mariae Virginis de Lourdes, anno MMXXII, Pontificatus Nostri
nono.} \quad Franciscus\endnote{Decree of Pope Francis of 11 February
2022, Latin text as published by the North American Province of the
Priestly Fraternity of Saint Peter (fssp.com/decree/), read 2026-07-25;
the same page carries an English rendering. The Fraternity's
international site publishes the English rendering alone, marked
``[Original: Latin and Spanish].'' Bounded negative: no publication of
this decree by the Holy See was located---not in \work{Acta Apostolicae
Sedis}, not on vatican.va, not in the Holy See Press Office bulletin---and
the canonical literature that treats the decree cites it from the
Fraternity. The English quoted below is the Fraternity's own; no project
translation of the Latin is offered.}
\end{quote}

The Fraternity's English renders this as: the Holy Father ``grants to
each and every member of the Society of Apostolic Life `Fraternity of
Saint Peter', founded on July 18, 1988 and declared of `Pontifical Right'
by the Holy See, the faculty to celebrate the sacrifice of the Mass, and
to carry out the sacraments and other sacred rites, as well as to fulfill
the Divine Office, according to the typical editions of the liturgical
books, namely the Missal, the Ritual, the Pontifical and the Roman
Breviary, in force in the year 1962''; they may use this faculty ``in
their own churches or oratories; otherwise it may only be used with the
consent of the Ordinary of the place, except for the celebration of
private Masses''; and ``Without prejudice to what has been said above,
the Holy Father suggests that, as far as possible, the provisions of the
motu proprio Traditionis Custodes be taken into account as well.''

\subsection{Six observations}

\emph{One: the decree is the decree of 10 September 1988, re-enacted by
the pope himself.} Compare the two operative sentences. The 1988 decree:
\latin{facultatem concedit Missae sacrificium celebrandi, ritus
sacramentorum aliosque sacros ritus peragendi necnon Officium Divinum
persolvendi secundum editiones typicas librorum liturgicorum anno 1962
vigentium, scilicet Missale, Rituale, Pontificale, et Breviarium
Romanum. Qua facultate uti poterunt in ecclesiis vel oratoriis propriis;
alibi vero nonnisi de consensu Ordinarii loci, excepta Missae privatae
celebratione.} The 2022 decree repeats the second sentence
word for word and the first with only the changes required by the new
grantor and the new grantees. Whatever else was happening in February
2022, the Holy See was not improvising: it was reissuing a thirty-four-year-old
concession over a papal signature.

\emph{Two: the grantees changed, and the change is a widening in one
direction and a narrowing in another.} The 1988 decree granted the
faculty to the Fraternity as a body (``concedes to that which is called
the Fraternity of St.\ Peter''), and the erection decree extended the use
of the books to members and to guests and ministers in the Fraternity's
churches. The 2022 decree grants it \latin{omnibus et singulis sodalibus},
to each and every member. That is a personal grant to every member
wherever he is---and it says nothing at all about guests, visiting priests,
or priests who minister in the Fraternity's churches without belonging to
it.

\emph{Three: the decree calls the Fraternity an institute of consecrated
life, and it gets its name wrong.} \latin{Instituti vitae consecratae
``Fraternitas Sancti Petri'' nuncupati}: the Fraternity is a society of
apostolic life, not an institute of consecrated life, and its name is
\latin{Fraternitas Sacerdotalis Sancti Petri}. The Fraternity's own
English rendering silently corrects the first error, writing ``the Society
of Apostolic Life `Fraternity of Saint Peter'.'' Neither error affects
the validity of a grant whose beneficiary is unmistakably identified by
its foundation date; both are worth recording, because a document that
misdescribes the canonical figure of its beneficiary is not the product of
a careful dicasterial process, and because the same slip had appeared in
the Commission's letter of 22 July 1988.

\emph{Four: the decree contains no reference to \work{Traditionis
custodes} art.~8, and does not derogate from it in terms.} It does not
say that the motu proprio does not bind the Fraternity. It does not say
that the 1988 concessions survive. It grants afresh---which is, in
practice, the most economical way of dealing with an abrogation, since
what is granted anew needs no defence against what abrogated the old
grant. But the technique leaves everything that was not re-granted
unprotected.

\emph{Five: the clause about \work{Traditionis custodes} is hortatory,
and the verb is the whole point.} \latin{Suadet}: he advises, urges,
counsels. Not \latin{praecipit}, not \latin{mandat}, not \latin{statuit}.
And what is advised is not compliance but that thought be given---\latin{ut
sedulo cogitetur}---and even that only \latin{quantum fieri potest}, as far
as can be done. A clause of that construction imposes no obligation
enforceable by anyone, and it is preceded by \latin{quibus rite servatis},
the foregoing having been duly observed: the grant is not conditioned on
the advice.

\emph{Six: the decree has no term, and no promulgation.} It states no
expiry and no review. It was not published in \work{Acta Apostolicae
Sedis} and has not been located in any Holy See publication. Absence of a
term is not a guarantee of perpetuity---a faculty granted by the supreme
legislator can be revoked by him---and absence of promulgation does not
invalidate an act addressed to determinate persons. But between them the
two absences mean that the Fraternity's principal liturgical title at this
account's cutoff is a document whose only public witness is the
beneficiary.

\subsection{The audience of 29 February 2024}

Two years later the Fraternity reports a further audience. On 29 February
2024 Francis received the superior general, Father Andrzej Komorowski,
accompanied by Father Benoît Paul-Joseph and Father Vincent Ribeton. In
the Fraternity's account the meeting was ``an opportunity for them to
express their deep gratitude to the Holy Father for the decree of
February 11, 2022, by which the Pope confirmed the liturgical specificity
of the Fraternity of St.\ Peter, but also to share with him the
difficulties encountered in its application.'' The pope, it says, ``was
very understanding and invited the Fraternity of St.\ Peter to continue to
build up ecclesial communion ever more fully through its own proper
charism.''\endnote{``Audience with Pope Francis'' (fssp.org), Fribourg, 1
March 2024, read 2026-07-25. Party account of a private audience. No Holy
See record of the audience was located.}

``Difficulties encountered in its application'' is the Fraternity's own
phrase and it is the closest thing to a public admission that the decree
of 2022 has not settled the matter on the ground. The account does not
say what the difficulties are; neither does any document reached here.
The obvious candidates are the ones the \work{Responsa} of 2021 created:
a bishop who reads the general law as governing what may happen in his
diocese, a parish church that may not be used, a Pontifical that a bishop
may not permit. But that is inference, and it is labelled as such.

\actrecord{Decree of Francis, 11 February 2022 (Latin, published by the
Fraternity).}{That the pope personally granted to each and every member of
the Fraternity the faculty to celebrate Mass, the sacraments and other
sacred rites and to say the Divine Office according to the typical
editions of the Missal, Ritual, Pontifical and Breviary in force in 1962;
in the Fraternity's own churches and oratories by the faculty itself,
elsewhere only with the local Ordinary's consent, private Masses
excepted; and that, this being duly observed, he also advises that thought
be given, as far as possible, to the provisions of \work{Traditionis
custodes}.}{A papal faculty of one page, never promulgated in the
official gazette and known only from the beneficiary's publication. It
grants no canonical status---the Fraternity already had one---derogates
from no article of \work{Traditionis custodes} in terms, names no term or
review, says nothing about guests or non-member ministers, and attaches
its only reference to the motu proprio to a verb of counsel.}
